In this part, we move to the fourth and final (for this book) step in the transducer ladder, namely the  polyregular functions. All transducer classes that we have seen before had linear output size, and the purpose of the polyregular functions is to achieve polynomial output size. 

As was the case for the regular functions, probably the quickest way to define the polyregular functions is in terms of prime functions, so we begin with that. We will simply add one more prime function, which has quadratic output size. 

\begin{myexample}[Marked squaring]\label{ex:marked-squaring}
    The marked squaring function is illustrated in the following example: 
    \begin{align*}
    1234 \quad \mapsto \quad \underline 1234 \ \underline{12} 34\  \underline{123} 4\  \underline{1234}.
    \end{align*}
    More formally, if the input  string has $n$ letters, then it is copied $n$ times, and in the $i$-th copy the  first $i$ letters are  underlined\footnote{The underlines might seem mysterious at first, but we will establish later in this book that they are necessary, and a version of the function without underlines would have smaller expresssive power.}. In our example, we have spaces in the output for readability, but the function does not produce them, i.e.~the output length is $n^2$. (Using a variant with spaces would not change the class of polyregular functions.) As was the case for map reverse and map duplicate, marked squaring is not one function, but a family of functions, with one function of type 
    \begin{align*}
    A^* \to (\myunderbrace{A+A}{two copies of the \\ input alphabet, one with \\ underlines and one without})^*
    \end{align*}
    for each input alphabet $A$.
\end{myexample}


\begin{definition}
    [Polyregular functions]\label{def:polyregular-functions}
    A string-to-string function is  \emph{polyregular} if it can be obtained as a finite composition of functions, each of which is either regular, or the marked squaring function from Example~\ref{ex:marked-squaring}. In other words: 
    \begin{align*}
    \text{polyregular} \quad \eqdef \quad \text{(regular $\cup$ marked squaring)}^*.
    \end{align*}
\end{definition}

Of course, the regular functions can be further decomposed into primes, using prime decompositions of regular functions, rational functions, and Mealy machines. The above definition is pleasingly short, but it does not adequatlely convey the expressive power and algorithmic strength of polyregular functions. To do that, we will present several characterisations of the polyregular functions in this part, which models coming from different fields, such as automata or programming languages. 

Before we do that, let us discuss our favourite three properties of transducer classes: continuity, closure under composition, and decidable equivalence. For closure under composition, there is nothing to do, since the class is defined in terms of composition. Decidable equivalence, regrettably, is not known and is a major open problem in this field. We are left with the continuity. 

\begin{theorem}\label{thm:polyregular-functions-are-continuous}
    Polyregular functions are continuous.
\end{theorem}
\begin{proof}
    Since continuous functions are closed under composition, and regular functions are continuous by Theorem~\ref{thm:regular-functions-are-continuous-and-closed-under-composition}, it remains to show that marked squaring is continuous. Consider then a composition 
    \begin{align}\label{eq:marked-squaring-composition}
    A^* \xrightarrow{\text{marked squaring}} (A+A)^* \xrightarrow L \Bool
    \end{align}
    of marked squaring with a regular language $L$. We want to show that this composition, which corresponds to the inverse image of $L$ under marked squaring, is a regular language. Suppose that the language $L$ is recognised by a \dfa  $\Aa$ with states $Q$.     For a string $w  \in A^*$ and a position $i$ in this string, let us write 
    \begin{align*}
    \delta_w^x : Q \to Q
    \end{align*}
    for the state transformation of the automaton $\Aa$ when reading the string that is obtained from $w$ by underlining the first $i$  input positions.  We will decompose the language in~\eqref{eq:marked-squaring-composition} into two steps, such that the first step is a rational function, and the second step is a regular language. Composing these two steps will give a regular language, by continuity of rational functions, and will hence prove that the language in~\eqref{eq:marked-squaring-composition} is regular. 
    \begin{enumerate}
        \item Suppose that the input is a string $w \in A^*$.  In the first step, we replace each input position $i$ by the state transformation $\delta_w^i$. The result is a string which has the same length as $w$, but is over the alphabet of state transformations $Q \to Q$. We will argue below that this first step is a rational function.

        \item If we compose the state transformations in the output of the first step, then we get the state transformation of the automaton $\Aa$ on the result of applying marked squaring to $w$. Therefore, in the second step, we: (a) compose all the state transformations in the output of the first step; and then  (b) check if the resulting state transformation takes the initial state of $\Aa$ to an accepting state. The second step can be implemented by a \dfa, whose states are state transformations, and is hence a regular language. 
    \end{enumerate}

As we have explained, the composition of the above two steps gives the same language as the composition in~\eqref{eq:marked-squaring-composition}, and hence proves regularity of the latter language. 

To conclude the proof, we will explain why the first step is a rational function. We will compute it using an unambiguous one-way transducer with output, in which the states are pairs $(\delta,\gamma)$ of state transformations. The invariant is going to be  that: (*) $\delta$ represents the underlined state transformation of the past, and $\gamma$ represents the non-underlined state transformation of the future. The transitions are of the form 
        \begin{align*}
        (\delta_1,\gamma_1) \xrightarrow {a / \delta_2 \cdot \gamma_2} (\delta_2, \gamma_2)
        \end{align*}
        such that the following conditions are satisfied:
        \begin{align*}
        \delta_2 = \delta_1 \cdot \myunderbrace{\delta_{\underline a}}{underlined state \\ transformation of $a \in A$} \qquad \text{and} \qquad  \gamma_1 = \myunderbrace{\delta_{a}}{non-underlined state \\ transformation of $a \in A$} \cdot \gamma_2.
        \end{align*}
        The initial states are those where $\delta$ is the identity transformation, and the accepting states are those where $\gamma$ is the identity transformation. This automaton has a unique accdepting run on every input string, which satisfies the invariant (*) described above.
\end{proof}
     